% Generated draft: requires mathematical and source review.
\documentclass[11pt]{article}
\usepackage[margin=25mm]{geometry}
\usepackage{fontspec}
\setmainfont{Latin Modern Roman}
\usepackage{amsmath,amssymb,mathtools,unicode-math}
\setmathfont{Latin Modern Math}
\usepackage{enumitem,needspace,xurl,hyperref}
\setlength{\parindent}{0pt}
\setlength{\parskip}{5pt}
\setlength{\emergencystretch}{3em}
\newcommand{\N}{\mathbb{N}}
\newcommand{\Z}{\mathbb{Z}}
\newcommand{\Q}{\mathbb{Q}}
\newcommand{\R}{\mathbb{R}}
\newcommand{\C}{\mathbb{C}}
\newcommand{\F}{\mathbb{F}}
\newcommand{\A}{\mathbb{A}}
\newcommand{\PP}{\mathbb{P}}
\newcommand{\E}{\mathbb{E}}
\newcommand{\eps}{\varepsilon}
\newcommand{\dd}{\,\mathrm{d}}
\begin{document}
\section*{UnsolvedMath ID: 30002653}
\textbf{Source problem code:} catalogue entry
\subsection{Definitions and mathematical statement}
Write \(P_g^{\mathrm{Perf}}\colon X_g\dashrightarrow Y_g^{\mathrm{Perf}}\) for the rational extension of the Prym map mentioned in the supplied record, where \(X_g\) denotes its intended compactified source and \(Y_g^{\mathrm{Perf}}\) its intended perfect-cone target. These symbols are placeholders: the record does not specify the source compactification, the target dimension, the base field, or the admissible values of \(g\).

For specified source and target spaces, let \(U_g\subseteq X_g\) be the locus of points having an open neighbourhood on which this rational map extends to a regular morphism. Its indeterminacy locus is
\[
\operatorname{Indet}(P_g^{\mathrm{Perf}})=X_g\setminus U_g.
\]
The symbols \(FS_2\) and \(FS_3\) refer to the two Friedman--Smith loci named in the supplied assessment. Their precise definitions, ambient spaces, and closure conventions are not supplied.

The question recorded by the source is whether
\[
\operatorname{Indet}(P_g^{\mathrm{Perf}})=FS_2\cup FS_3.
\]
This formula preserves the stated target question, but a fully specified mathematical assertion requires definitions of the moduli spaces and Friedman--Smith loci, a range and quantifier for \(g\), and clarification of whether equality concerns underlying sets, reduced closed loci, or scheme structures. The supplied assessment mentions higher genus without specifying a numerical range.
\subsection{Short English statement}
Does the rational extension of the Prym map to the perfect-cone compactification fail to be regular precisely along the union of the Friedman--Smith loci \(FS_2\) and \(FS_3\)? The source leaves the moduli conventions, genus range, and precise meanings of these loci unresolved.
\subsection{Sources}
\begin{itemize}
\item Local registry: \texttt{data/unsolvedmath/problems.json}, permanent UnsolvedMath ID \(30002653\), source problem code \texttt{catalogue entry}; collection: Oberwolfach Reports Open Problems. No individual proposer is supplied. The proposed year is recorded as 2014.
\item Supplied report citation: \emph{Komplexe Analysis} (2015), Oberwolfach Reports, pp.~2177--2234, DOI \texttt{10.4171/owr/2014/39}. \url{https://doi.org/10.4171/owr/2014/39}. This is also the supplied literature link labelled ``original Oberwolfach report.'' The citation's 2015 year and the record's 2014 attribution are retained without reconciliation.
\item Supplied literature reference: Casalaina-Martin et al., ``Prym compactifications,'' DOI \texttt{10.4171/JEMS/678}. \url{https://doi.org/10.4171/JEMS/678}. No more detailed bibliographic information or pinpoint reference is supplied.
\item Literature assessment dated exactly \texttt{2026-08-21}: the record labels the problem open and states that established inclusions bound the perfect-cone Prym indeterminacy locus in general, while equality with the two Friedman--Smith loci remains unresolved in higher genus. This is supplied dated curation derived from the OWR material, not an independently verified literature assessment; the status and cited claims require independent verification.
\item The registry labels the statement \texttt{corrected\_verified} and reports that notation, typography, and wording were normalized without changing the mathematical content. These are supplied editorial claims, not independent mathematical verification of this draft.
\end{itemize}
\end{document}
